%----------------------------------------------------------------------------------------
% Tailored CV — Wiss. Mitarbeiter*in, AG Neurobiologie, Institut für Biologie
% Freie Universität Berlin. Kennung 21259000_2026. Deadline 29.06.2026. E13 TV-L.
% HONEST computational framing: image analysis + connectome + transcriptome + comp. neuro.
% Does NOT claim Drosophila, 2-photon/confocal facility operation, or a completed PhD.
%----------------------------------------------------------------------------------------

\documentclass[10pt,a4paper,sans]{moderncv}

\usepackage[utf8]{inputenc}
\usepackage[T1]{fontenc}
\usepackage{lmodern}

\moderncvstyle{classic}
\moderncvcolor{green}

\usepackage[scale=0.78]{geometry}
\usepackage{ragged2e}

%----------------------------------------------------------------------------------------

\firstname{Kundai Farai}
\familyname{Sachikonye}

\address{Zeltnerstrasse 30}{90443, Nuernberg, Germany}
\mobile{(+49) 1626386344}
\email{kundai.sachikonye@bitspark.com}
\homepage{github.com/fullscreen-triangle}
\photo[70pt][0.4pt]{pictures/Bew2.jpg}

\quote{Computational scientist: digital image analysis of microscopy data,
neural-circuit / connectome modelling, and transcriptome analysis.}

%----------------------------------------------------------------------------------------

\begin{document}

\makecvtitle

%----------------------------------------------------------------------------------------
%	PROFILE
%----------------------------------------------------------------------------------------

\section{Profile}

\cvitem{}{
  Computational researcher focused on \textbf{quantitative analysis of imaging
  data}, \textbf{neural-circuit / connectome modelling}, and \textbf{transcriptome
  analysis}. My background is computational rather than wet-lab Drosophila
  genetics or laser-microscope operation; I would contribute to the AG
  Neurobiologie's complementary computational and image-analysis work.
}

%----------------------------------------------------------------------------------------
%	SELECTED COMPUTATIONAL WORK
%----------------------------------------------------------------------------------------

\section{Selected Computational Work}

\cvitem{Microscopy image analysis}{
  \textbf{Quantitative image-analysis framework} --- segmentation, PSF
  deconvolution, scale-field / coordinate reconstruction, and morphometry with
  uncertainty quantification, for fluorescence and phase-contrast microscopy.
  GPU-accelerated; live tool.
  \href{https://hieronymus-omega.vercel.app/}{hieronymus-omega.vercel.app}
}

\cvitem{Computational neuroscience}{
  \textbf{Neural-circuit modelling} (project "Olduvai") --- closed-circuit
  decomposition of motor/neural signals (supraspinal vs.\ peripheral),
  applied to neurodegeneration (Parkinson's, ataxia) from physiological
  time-series; circuit-graph (connectome-style) network analysis.
}

\cvitem{Transcriptome / multi-omics}{
  \textbf{gospel} --- integration of transcriptomics with genotype and
  metagenomics; validated against 1000~Genomes, gnomAD, UK~Biobank.
  \href{https://github.com/fullscreen-triangle/gospel}{github.com/fullscreen-triangle/gospel}
}

%----------------------------------------------------------------------------------------
%	EXPERIENCE
%----------------------------------------------------------------------------------------

\section{Experience}

\cventry{2021--present}{Independent Computational Researcher}{Bitspark GmbH}{}{}{
  Full-time computational research: quantitative image analysis,
  neural-circuit / connectome modelling, transcriptome and multi-omics analysis.
  Scripted, reproducible pipelines with statistical evaluation; open, documented,
  deployed tools.
}

\cventry{2019--2021}{Computational Lipidomics --- Bioinformatics, HPC \& Teaching}{Technische Universität München}{Munich}{}{
  High-throughput data analysis at HPC scale; statistical evaluation of
  large datasets. Taught at Master's level and supervised student work.
}

\cventry{2017--2018}{Glycomics and Proteomics}{Universitätsklinikum Hamburg-Eppendorf}{Hamburg}{}{
  LC-MS/MS and MALDI-TOF data pipelines; reproducible high-throughput workflows.
}

\cventry{2013}{Cancer Biomarker Discovery}{University Medical Center Groningen}{Groningen}{}{
  UPLC-MS/MS proteomics; statistical analysis of high-dimensional data.
}

%----------------------------------------------------------------------------------------
%	EDUCATION
%----------------------------------------------------------------------------------------

\section{Education}

\cventry{2019--2021}{Doctoral research, Computational Lipidomics}{Technische Universität München}{}{}{
  (Doctoral research; not completed.) Computational analysis of
  high-throughput data.
}
\cventry{2018--2019}{MSc Computational Biology}{Universität Tübingen}{}{}{
  Sequence/transcriptome analysis, statistics, ML for biological data.
}
\cventry{2014--2017}{MSc Pharmaceutical Biotechnology}{HAW Hamburg}{}{}{}
\cventry{2011--2014}{BSc Biochemistry and Cell Biology}{Jacobs University Bremen}{}{}{
  Thesis: DNA interstrand crosslinking --- mammalian cell culture and
  fluorescence microscopy.
}

%----------------------------------------------------------------------------------------
%	TECHNICAL SKILLS
%----------------------------------------------------------------------------------------

\section{Technical Skills}

\cvitem{Image analysis}{
  Digital image analysis (segmentation, deconvolution, registration,
  morphometry); statistical evaluation of imaging data; GPU/GLSL image processing
}

\cvitem{Computational neuroscience}{
  Neural-circuit and connectome-style network/graph analysis; physiological
  time-series modelling
}

\cvitem{Data analysis}{
  Transcriptome analysis; statistics; scripted, reproducible pipelines
  (Python, R); HPC (Ray, Dask, SLURM)
}

\cvitem{Languages (programming)}{
  Python (primary), Rust, R, Bash, SQL
}

%----------------------------------------------------------------------------------------
%	LANGUAGES & TEACHING
%----------------------------------------------------------------------------------------

\section{Teaching \& Languages}

\cvitem{Teaching}{
  Taught at Master's level at TUM (2019--2021); supervised student work.
}
\cvitemwithcomment{English}{Native / Fluent (C2)}{lab language}
\cvitemwithcomment{German}{Conversational}{resident in Germany since 2011}

%----------------------------------------------------------------------------------------

\renewcommand{\listitemsymbol}{-~}

\end{document}
